%==========================================================================
%  TRES0312 Fysikk — Foiler-mal
%
%  MØrk/lys modus: kommenter ut én av de to linjene nedenfor.
%==========================================================================
\newif\ifdarkmode
\darkmodetrue      % ← mørk modus
%\darkmodefalse   % ← lys modus  (fjern % foran og legg % foran linjen over)

\documentclass[aspectratio=169, 12pt]{beamer}

\usepackage{amd-beamer}
\usetikzlibrary{overlay-beamer-styles}   % for "visible on=<->" i enkelte tikz-figurer

%--------------------------------------------------------------------------
%  DOKUMENTINFO
%--------------------------------------------------------------------------
\title{Kinematikk}
\subtitle{TRES0312 --- Plenumsregning}
\author{Ben David Normann}
\date{\today}

%--------------------------------------------------------------------------
%  SEKSJONSKONTEKST
%  \setcounter{section}{N} setter gjeldende seksjonsnummer til N, slik at
%  \defn, \exm og \oppgave får samme nummerering som i kompendiet.
%  Tellerne nullstilles automatisk ved hvert \setcounter{section}{...}.
%--------------------------------------------------------------------------
\setcounter{section}{3}   % Kapittel 3: "Kinematikk"

\begin{document}

%--------------------------------------------------------------------------
%  SLIDE 1 — TITTELSIDE
%--------------------------------------------------------------------------
\begin{frame}[plain]
  \titlepage
\end{frame}

%--------------------------------------------------------------------------
%  SLIDE 2 — Plan for økta
%--------------------------------------------------------------------------
\begin{frame}{Plenumsregning}

  \oppsum{I dag}{%
  Vi går sammen gjennom flere oppgaver i kinematikk --- to i 1D og to i 2D --- som repetisjon før vi går videre til sirkelbevegelse.
  \begin{itemize}
    \item[1D] Bremsing, og fritt fall
    \item[2D] Kast fra bakkenivå, og kast fra en høyde
  \end{itemize}
  }

\end{frame}

%--------------------------------------------------------------------------
%  SLIDE 3 — Aktivitet: Bremsing (1D)
%--------------------------------------------------------------------------
\begin{frame}{Aktivitet (1D)}

  \aktivitet{Bremsing}{%
  En bil kjører med fart $v_0 = 20\,\mathrm{m/s}$ og bremser med konstant akselerasjon $a = -5\,\mathrm{m/s^2}$ helt til den stopper.
  \begin{enumerate}[a)]
    \item Hvor lang tid bruker bilen på å stoppe?
    \item Hvor langt kjører bilen i løpet av oppbremsingen?
  \end{enumerate}
  }

\end{frame}

%--------------------------------------------------------------------------
%  SLIDE 4 — Aktivitet: Fritt fall fra bro (1D)
%--------------------------------------------------------------------------
\begin{frame}{Aktivitet (1D)}

  \aktivitet{Fritt fall fra en bro}{%
  En stein slippes (uten starthastighet) fra en bro $20\,\mathrm{m}$ over vannflaten.
  \begin{enumerate}[a)]
    \item Hvor lang tid tar det før steinen treffer vannet?
    \item Hvor fort beveger steinen seg i det den treffer vannet?
  \end{enumerate}
  }

\end{frame}

%--------------------------------------------------------------------------
%  SLIDE 5 — Aktivitet: Ball kastet fra bakkenivå (2D)
%--------------------------------------------------------------------------
\begin{frame}{Aktivitet (2D)}

  \aktivitet{Ball kastet fra bakkenivå}{%
  En ball kastes med utgangsfart $v_0 = 15\,\mathrm{m/s}$ i en vinkel på $40^\circ$ over et flatt, horisontalt underlag ($y_0 = 0$).
  \begin{enumerate}[a)]
    \item Finn $v_0$ dekomponert i $x$- og $y$-retning.
    \item Bestem hvor lang tid ballen er i lufta.
    \item Bestem hvor langt unna ballen lander, og hvor høyt den når på det høyeste.
  \end{enumerate}
  }

\end{frame}

%--------------------------------------------------------------------------
%  SLIDE 6 — Aktivitet: Ball ruller av et bord (2D)
%--------------------------------------------------------------------------
\begin{frame}{Aktivitet (2D)}

  \aktivitet{Ball ruller av et bord}{%
  En ball ruller av kanten på et bord med høyde $h = 1.20\,\mathrm{m}$, med horisontal fart $v_0 = 2.5\,\mathrm{m/s}$ i det den forlater bordkanten.
  \begin{enumerate}[a)]
    \item Hvor lang tid bruker ballen på å falle ned til gulvet?
    \item Hvor langt fra bordkanten treffer ballen gulvet?
    \item Hva er farten (størrelsen) til ballen i det den treffer gulvet?
  \end{enumerate}
  }

\end{frame}

%--------------------------------------------------------------------------
%  SLIDE 7 — Ekstra: for de som blir tidlig ferdig
%--------------------------------------------------------------------------
\begin{frame}{Ekstra (for de som blir tidlig ferdig)}

  \aktivitet{To biler møtes (1D)}{%
  Bil A starter i ro ved $x=0$ og akselererer med $a = 2\,\mathrm{m/s^2}$. Samtidig kjører bil B med konstant fart $v = 15\,\mathrm{m/s}$ fra $x = 300\,\mathrm{m}$, i retning mot bil A.
  Når og hvor møtes bilene?
  }

  \pause
  \aktivitet{Finn utgangsfarten (2D)}{%
  En pil skytes i en vinkel på $45^\circ$ over et flatt, horisontalt underlag, og lander $30\,\mathrm{m}$ unna --- på samme høyde den ble skutt fra ($y_0 = 0$).
  Hva var utgangsfarten til pilen?
  }

\end{frame}

%--------------------------------------------------------------------------
%  SLIDE 8 — Neste gang
%--------------------------------------------------------------------------
\begin{frame}
\begin{center}
  \Huge{Neste gang: Sirkelbevegelse}
  \begin{itemize}
    \item Sentripetalakselerasjon
    \item Sirkelbevegelse og omløpstid
  \end{itemize}
\end{center}
\end{frame}

\end{document}
